\begin{proof}[Proof of \cref{obj:thm:mixed-count-upper}]
\begin{enumerate}
\item Choose the universal constant \(\overline C>0\) supplied by
\cref{obj:thm:uniform-mixed-count-certificate}. Fix \(n,d,q\) and a full-data law \(P\) satisfying the stated conditions. Membership in \(\mathcal M_{n,d,q}\) supplies all hypotheses of that certificate.
By \cref{obj:lem:cell-identification-background}, the causal target equals the total observed-cell functional on this model class, which is the identification step used by the certificate. The certificate therefore gives
\[
E_P\!\left[
\left(
\widehat\tau^{\mathrm{mix}}_{n,d,q}(\mathbf O_n)-\tau(P)
\right)^2
\right]
\le \mathfrak U_{n,d,q}
\le \overline C\,r_{n,d,q}.
\]
The envelope and branch rules in that certificate are exactly those in the present statement. It remains to establish the comparison with the auxiliary randomized risk.
% lean: CausalSmith.Stat.MarRareqLogfrontier.uniform_mixed_count_certificate

\item Write the prefix probabilities and their remaining mass as
\[
w_{k_0}:=\Pr(K_0=k_0)
=e^{-n/2}\frac{(n/2)^{k_0}}{k_0!},
\qquad k_0\in\mathbb N,
\qquad
w_{\mathrm{tail}}:=1-\sum_{k_0=0}^{n}w_{k_0}.
\]
The singleton prefix events are disjoint, so
\[
w_{k_0}\ge0,
\qquad
\sum_{k_0=0}^{n}w_{k_0}\le1,
\qquad
w_{\mathrm{tail}}=\Pr(K_0>n)\ge0.
\]
For each integer \(0\le k_0\le n\), define the assignment set by
\[
\mathcal A_{\mathrm{str},k_0}
:=
\left\{
\sigma:\{1,\ldots,k_0\}
\longrightarrow
\{\mathrm{membership},\mathrm{pilot},\mathrm{estimation}\}
\right\},
\qquad
|\mathcal A_{\mathrm{str},k_0}|=3^{k_0}.
\]
Let the output at a prescribed prefix and assignment, and its assignment average, be
\[
\widetilde\tau_{k_0,\sigma}(\mathbf O_n)
:=
\widetilde\tau(\mathbf O_n;\omega)
\quad\text{when }\omega\text{ prescribes }K_0=k_0
\text{ and assignment }\sigma,
\]
\[
\overline\tau_{\mathrm{str},k_0}(\mathbf O_n)
:=
\frac{1}{3^{k_0}}
\sum_{\sigma\in\mathcal A_{\mathrm{str},k_0}}
\widetilde\tau_{k_0,\sigma}(\mathbf O_n).
\]
These outputs use the branch rules of
\cref{obj:def:mixed-count-estimator}. For \(k_0=0\), the assignment set contains the unique empty assignment.

The zero output on the prefix tail and the conditional-average definition give
\[
\widehat\tau^{\mathrm{mix}}_{n,d,q}(\mathbf O_n)
=
\sum_{k_0=0}^{n}
w_{k_0}\overline\tau_{\mathrm{str},k_0}(\mathbf O_n).
\]
Pointwise in the observed sample, its centered value therefore satisfies
\[
\widehat\tau^{\mathrm{mix}}_{n,d,q}(\mathbf O_n)-\tau(P)
=
w_{\mathrm{tail}}\bigl(-\tau(P)\bigr)
+
\sum_{k_0=0}^{n}
w_{k_0}
\bigl(\overline\tau_{\mathrm{str},k_0}(\mathbf O_n)-\tau(P)\bigr).
\]
The weights on the right are nonnegative and sum to one. Convexity of the square yields
\[
\left(
\widehat\tau^{\mathrm{mix}}_{n,d,q}(\mathbf O_n)-\tau(P)
\right)^2
\le
w_{\mathrm{tail}}\tau(P)^2
+
\sum_{k_0=0}^{n}
w_{k_0}
\left(
\overline\tau_{\mathrm{str},k_0}(\mathbf O_n)-\tau(P)
\right)^2.
\]
% lean: CausalSmith.Stat.MarRareqLogfrontier.finite_weighted_square_with_remainder

\item For every \(0\le k_0\le n\), centering the assignment average gives
\[
\overline\tau_{\mathrm{str},k_0}(\mathbf O_n)-\tau(P)
=
\frac{1}{3^{k_0}}
\sum_{\sigma\in\mathcal A_{\mathrm{str},k_0}}
\left(
\widetilde\tau_{k_0,\sigma}(\mathbf O_n)-\tau(P)
\right).
\]
The finite Cauchy--Schwarz inequality consequently gives
\[
\left(
\overline\tau_{\mathrm{str},k_0}(\mathbf O_n)-\tau(P)
\right)^2
\le
\frac{1}{3^{k_0}}
\sum_{\sigma\in\mathcal A_{\mathrm{str},k_0}}
\left(
\widetilde\tau_{k_0,\sigma}(\mathbf O_n)-\tau(P)
\right)^2.
\]
Multiplying these inequalities by the nonnegative \(w_{k_0}\), summing, and inserting the result into the preceding bound yields
\[
\begin{aligned}
&\left(
\widehat\tau^{\mathrm{mix}}_{n,d,q}(\mathbf O_n)-\tau(P)
\right)^2\\
&\quad\le
w_{\mathrm{tail}}\tau(P)^2
+
\sum_{k_0=0}^{n}\frac{w_{k_0}}{3^{k_0}}
\sum_{\sigma\in\mathcal A_{\mathrm{str},k_0}}
\left(
\widetilde\tau_{k_0,\sigma}(\mathbf O_n)-\tau(P)
\right)^2.
\end{aligned}
\]
% lean: CausalSmith.Stat.MarRareqLogfrontier.finite_uniform_square

\item Conditional on the observed sample, the prefix has probabilities \(w_{k_0}\), each assignment has probability \(3^{-k_0}\), and the tail output is zero. Thus
\[
\begin{aligned}
&E_\omega\!\left[
\left(\widetilde\tau(\mathbf O_n;\omega)-\tau(P)\right)^2
\mid\mathbf O_n
\right]\\
&\quad=
w_{\mathrm{tail}}\tau(P)^2
+
\sum_{k_0=0}^{n}\frac{w_{k_0}}{3^{k_0}}
\sum_{\sigma\in\mathcal A_{\mathrm{str},k_0}}
\left(
\widetilde\tau_{k_0,\sigma}(\mathbf O_n)-\tau(P)
\right)^2.
\end{aligned}
\]
The observed-sample space is finite, so the displayed functions are measurable and integrable. Integrating the pointwise inequality establishes
\[
E_P\!\left[
\left(
\widehat\tau^{\mathrm{mix}}_{n,d,q}(\mathbf O_n)-\tau(P)
\right)^2
\right]
\le
E_{P,\omega}\!\left[
\left(
\widetilde\tau(\mathbf O_n;\omega)-\tau(P)
\right)^2
\right].
\]
Together with the certificate bound, this proves all the asserted inequalities with the same universal constant \(\overline C\).
% lean: CausalSmith.Stat.MarRareqLogfrontier.mixed_count_derandomization
% lean: CausalSmith.Stat.MarRareqLogfrontier.mixed_count_upper
\end{enumerate}
\end{proof}
